% TT09031C
% Rückenmarksverletzung, Verdacht
% 14.0
% 2013-11-30
\label{m:ws-trauma-verdacht}
\EE{Taktik}: \Speech{
Immobilisation und Transport in eine geeignete Einrichtung
zur Abklärung}


Vgl. \prettyref{m:einschaetzung-indikation-ws-immobilisation}
\begin{itemize}
  \item \EE{\Abk{HWS}-Immobilisation} (HWS-Schiene)
  \item Bergung mit Schaufeltrage, Spineboard oder
  KED-System
  \item \EE{Lagerung} und Schienung mittels Vakuummatratze, Schaufeltrage mit
  Fixationsmöglichkeit oder Spineboard.\footnote{Welches Hilfsmittel
  eingesetzt wird ist abhängig vom Patientenzustand,
  dem Training des Teams, internen Vorschriften und
  vorhandenem Material.}
  \item \EE{Monitoring}: Achte besonders auf Störungen
  von Atmung und Bewusstsein, Schockentwicklung und
  beginnende neurologische Ausfälle.
%   \item \EE{Unfallmechanismus} erheben und dokumentieren
  \item \EE{Transportentscheidung}: Je nach
  lokalen Gegebenheiten
\end{itemize}